\section{Penelitian Terkait}
Pada sub bab ini, dipaparkan mengenai penelitian-penelitian yang berkaitan dengan topik penelitian.

\begin{enumerate}
    \item Penelitian yang dilakukan oleh \citep{fuad2021covid19} yaitu memodelkan dan memprediksi penyebaran kasus COVID-19 di tia kota wilayah Mekkah yaitu Jeddah, Makkah dan Taif dengan memanfaatkan pendeketana statistic ARIMA dan STARIMA (Space Time ARIMA).
    Penelitian ini bertujuan untuk membandingkan performa kedua model tersebeut serta menilai pengaruh hubungan spasial antar wilayah terhadap Akurasi Prediksi.
    Dalam studinya, penulis menggunakan data 31 Mei hingga 11 Oktober 2020.
    Analisis awal menunjukkan bahwa data belum stasioner sehingga dilakukan transformasi log dan diferensiasi pertama.
    Model ARIMA dibangun secara terpisah untuk setiap kota Menggunakan pendekatan Box Jenkins, sedangkan model STARIMA dikembangkan secara simultan dengan mempertimbangkan hubungan spasial antar lokasi melalui spatial weigt matrix, yang merepresentasikan keterhubungan antar kota. \\
    Hasil penelitian menunjukkan bahwa model STARIMA mampu menangkap autokorelasi ruang waktu dengan lebih baik dibandingkan ARIMA univariat.
    Model terbaik yang diperoleh adalah STARIMA $(1_{0},1,1_{1})$ yang menunjukkan nilai MAE, MSE, dan RMSE lebih rendah dibandingkan model ARIMA diketioga kota.
    Temuan ini menegaskan bahwa mempertimbangkan hubungan spasial (misalnya kedekatan geografis atau mobilitas penduduk antar kota) dapat secara signifikan meningkatkan Akurasi Peramalan penyebaran COVID-19. \\
    Penelitian ini menyimpulkan bahwa pendekatan spatial temporal seperti STARIMA lebih unggul dan direkomendasikan untuk data epidemiologi yang memiliki pola penyebaran antar wilayah.
    Studi ini juga membuka peluang untuk penggunaan model spatial temporal yang lebih kompleks seperti Bayesian spatial temporal models atau spatial panel models untuk analisis lanjutan.

    \item Dalam penelitian yang berjudul “\textit{Enhanced Spatio-Temporal Modeling for Rainfall Forecasting: A High-Resolution Grid Analysis}” Curah hujan sangat penting bagi pertanian, terutama di negara seperti India. Karena bersifat spasio-temporal, peramalan curah hujan menjadi kompleks dan tidak cocok jika hanya menggunakan model deret waktu atau machine learning konvensional. Diperlukan pendekatan yang mampu menangani pola \textit{non-linear, non-stasioner, dan non-normal}. Model STARMA efektif dalam menangkap ketergantungan spasial dan temporal sekaligus, dibantu oleh Spatial Weight Matrix (SWM) untuk mengintegrasikan pengaruh lokasi sekitar. Studi ini menggunakan data curah hujan tahunan selama 50 tahun (1970–2019) dari 119 lokasi di West Bengal, India, yang dibagi menjadi data pelatihan dan pengujian. Hasilnya menunjukkan bahwa model STARMA mampu meningkatkan akurasi peramalan, yang bermanfaat untuk perencanaan dan pengelolaan pertanian di wilayah rentan terhadap perubahan iklim \citep{alam2024enhanced}
    \item Dalam penelitian yang berjudul\textit{ “An improved Space-Time Autoregressive Moving Average (STARMA) model for Modelling and Forecasting of Spatio-Temporal time-series data”} menjelaskan bahwa Model\textit{ Box-Jenkins } univariat tidak mampu menangani ketergantungan spasial yang sering muncul dalam fenomena iklim. Untuk mengatasi hal ini, digunakan model STARMA\textit{ (Space-Time Autoregressive Moving Average)} yang memasukkan hubungan spasial dan temporal secara simultan melalui matriks bobot spasial. Dalam penelitian ini, digunakan matriks bobot spasial seragam orde dua dan bobot berdasarkan jarak terbalik yang dihitung menggunakan jarak lingkaran besar \textit{(great circle distance)} dari koordinat lintang dan bujur. Metode ini diuji pada data simulasi dan data nyata berupa suhu maksimum bulanan di sembilan distrik Karnataka Utara. Hasilnya menunjukkan bahwa model STARMA yang diusulkan memberikan performa yang lebih baik dibandingkan ARIMA univariat dan STARMA orde spasial pertama dalam pemodelan dan peramalan deret waktu spasio-temporal \citep{rathod2018improved}
    \item Penelitian ini berhasil mengembangkan dan menguji model prediksi hasil panen padi berbasis pendekatan two-stage yang dinamakan STARMA II \textit{(Spatio-Temporal Autoregressive Moving Average with Time Delay Neural Network)}. Model ini menggabungkan keunggulan model STARMA yang menangkap ketergantungan spasial dan temporal secara linear dengan kemampuan model\textit{ Time Delay Neural Network } (TDNN) dalam mempelajari hubungan non-linear dari data time series. 
    
    Dalam implementasinya, pengaruh spasial antar distrik dimodelkan menggunakan \textit{Spatial Weight Matrix} (SWM) yang dibangun berdasarkan gabungan metode\textit{ K-Nearest Neighbors} (KNN) dan koefisien korelasi Pearson. Matriks ini berfungsi untuk merepresentasikan kemiripan spatio-temporal antar wilayah secara kuantitatif, yang kemudian digunakan dalam tahap pelatihan model untuk menangkap pola hubungan antar wilayah dalam konteks produksi padi. 
    
    Hasil evaluasi model menunjukkan bahwa STARMA II memberikan akurasi prediksi yang lebih tinggi dan konsisten dibandingkan dengan: 
    \begin{itemize}
        \item Model ARIMA yang hanya mempertimbangkan aspek temporal (waktu), dan 
        \item Model STARMA klasik yang hanya menangkap hubungan linier antar ruang dan waktu. 
    \end{itemize}
    
    Kinerja model diukur melalui metrik evaluasi seperti \textit{Root Mean Square Error} (RMSE), \textit{Mean Absolute Error} (MAE), dan \textit{Coefficient of Determination} (R²), di mana STARMA II unggul secara signifikan dalam semua indikator. 
    
    Dengan demikian, pendekatan two-stage hybrid model seperti STARMA II terbukti lebih efektif dalam memodelkan dinamika kompleks produksi padi yang dipengaruhi oleh faktor spasial dan temporal, serta memiliki potensi besar untuk diterapkan dalam sistem peringatan dini, perencanaan pertanian, dan pengambilan kebijakan berbasis data \citep{badu-apraku2021two}
    \item Penelitian yang dilakukan oleh Saha, Singh, Ray, Rathod, dan Dhyani (2022) bertujuan untuk meningkatkan akurasi peramalan suhu dengan mengembangkan model baru yang disebut \textit{Fuzzy Rule–Based Weighted } \textit{Space–Time Autoregressive Moving Average} (STARMA). Model ini merupakan pengembangan dari STARMA konvensional, dengan mengintegrasikan sistem inferensi fuzzy (FIS) untuk menghitung bobot spasial antar lokasi. Pendekatan ini dirancang untuk mengatasi keterbatasan dari matriks bobot spasial konvensional yang tidak mampu menangkap dinamika dan heterogenitas spasial secara optimal. Dalam penelitian ini, data suhu maksimum bulanan digunakan dari tujuh distrik di negara bagian West Bengal, India, untuk mengembangkan model, serta dari sembilan distrik di Karnataka untuk validasi. Hasil penelitian menunjukkan bahwa model fuzzy STARMA memberikan hasil yang lebih akurat dibandingkan dengan model STARMA konvensional dan model ARIMA. Hal ini dibuktikan dengan nilai \textit{Mean Absolute Percentage Error} (MAPE) yang lebih rendah, yaitu 1,78\% untuk model fuzzy STARMA, dibandingkan dengan 2,09\% untuk STARMA biasa, dan 4,39\% untuk ARIMA. Dengan demikian, penelitian ini menyimpulkan bahwa penggunaan sistem inferensi fuzzy dalam peramalan spasial-temporal memberikan peningkatan performa yang signifikan. Model ini sangat berpotensi untuk digunakan dalam aplikasi yang berkaitan dengan perencanaan iklim dan sektor pertanian, khususnya di wilayah yang memiliki karakteristik spasial yang kompleks seperti India \citep{saha2022fuzzy} 
    \item Penelitian yang dilakukan oleh Kumar, Sarkar, Dhakre, dan Bhattacharya (2023) bertujuan untuk meningkatkan akurasi peramalan suhu bulanan dengan mengembangkan pendekatan pemodelan hibrida yang menggabungkan model\textit{ Space–Time Autoregressive Moving Average } (STARMA) dan \textit{Generalized Autoregressive Conditional Heteroskedasticity } (GARCH). Studi ini difokuskan pada peramalan suhu maksimum bulanan dan rentang suhu di delapan lokasi di negara bagian Bihar, India. Model STARMA digunakan untuk menangkap pola spasial dan temporal dalam data, dengan matriks bobot spasial yang dikonstruksi berdasarkan jarak antar lokasi. Setelah penerapan STARMA, residual model dianalisis lebih lanjut menggunakan uji BDS dan ARCH-LM, yang menunjukkan adanya karakteristik nonlinier dan heteroskedastisitas. Untuk menangani fenomena ini, model GARCH diterapkan pada residual STARMA, membentuk pendekatan hibrida STARMA-GARCH. Hasil penelitian menunjukkan bahwa model STARMA-GARCH (khususnya konfigurasi STARMA(1,1,0,0)-GARCH(0,1)) memberikan performa yang lebih baik dibandingkan model STARMA murni, dengan nilai kesalahan prediksi (MAPE) yang lebih rendah di semua lokasi. Uji diagnostik tambahan seperti uji Ljung-Box Q juga menunjukkan bahwa residual dari model hibrida tidak memiliki autokorelasi, menandakan bahwa model tersebut telah dikalibrasi dengan baik. Kesimpulan dari penelitian ini adalah bahwa pendekatan pemodelan hibrida ini sangat efektif dalam menangkap dinamika kompleks suhu bulanan dan dapat menjadi alat yang andal untuk mendukung pengambilan keputusan di sektor-sektor yang sensitif terhadap iklim seperti pertanian dan pengelolaan sumber daya alam \citep{kumar_nodate}

    Berbagai penelitian sebelumnya telah banyak mengembangkan model peramalan spasio-temporal menggunakan pendekatan STARMA dan variannya dengan integrasi matriks bobot spasial. Meskipun model STARMA efektif dalam menangkap ketergantungan linier spasial dan temporal, model ini masih memiliki keterbatasan dalam mengakomodasi pola non-linear dan heterogenitas spasial yang kompleks, terutama dalam konteks data curah hujan. Selain itu, konstruksi matriks bobot spasial yang digunakan sebagian besar masih bersifat statis dan umum, seperti berdasarkan jarak geografis atau metode \textit{K-Nearest Neighbors}, sehingga kurang mampu menangkap dinamika hubungan spasial-temporal secara adaptif. Selanjutnya, belum banyak penelitian yang secara eksplisit memanfaatkan pembobotan deret waktu berbasis kemiripan spasio-temporal yang menggabungkan aspek spasial dan temporal secara simultan dan kuantitatif, khususnya untuk peramalan curah hujan dengan skala spasial dan temporal yang rinci. Selain itu, kompleksitas model yang tinggi pada beberapa pendekatan hibrida seperti \textit{neural network} dan \textit{fuzzy inference} juga menimbulkan tantangan terkait efisiensi komputasi pada data berukuran besar. Oleh karena itu, penelitian ini berupaya mengisi gap tersebut dengan mengembangkan model peramalan curah hujan yang menggunakan pembobotan seragam, pembobtan invers jarak, dan pembobotan korelasi silang yang adaptif. Pendekatan ini diharapkan dapat meningkatkan akurasi peramalan curah hujan sekaligus mendukung pengambilan keputusan di bidang pertanian dan mitigasi risiko bencana iklim.
\end{enumerate}